\uuid{g7q7}
\titre{Étude de suite}
\auteur{bodin ; hanani ; mzari}
\organisation{exo7}
\licence{CC BY-NC-SA 4.0}
\datecreate{2020-01-20}
\niveau{L1}
\module{Analyse}
\chapitre{Suite}
\sousChapitre{Autre}
\difficulte{3}
\theme{Suites, étude de suite}
\exo7qcmid{lille-1.suites.028}

\contenu{
\begin{question}
\qclassification{121.00}
\qtags{motcle=étude de suite}
Soit $(u_n)$ une suite croissante. On suppose que $\displaystyle 0\leq u_{n+1}\leq \frac{1}{2}u_n+\frac{1}{2^n}$ pour tout $n\geq 0$. Quelles sont les bonnes réponses ?
\begin{answers}
    \good{$(u_n)$ est majorée.}
    \bad{$(u_n)$ est divergente.}
    \good{$(u_n)$ est convergente et $\displaystyle 0\leq \lim _{n\to +\infty}u_n\leq 2$.}
    \good{$u_n=0$ pour tout $n\geq 1$.}
\end{answers}
\begin{explanations}
On vérifie par récurrence que, pour tout $n\geq 1$,
$$\displaystyle 0\leq u_n\leq \frac{1}{2^n}u_0+1+\frac{1}{2}+\frac{1}{2^2}+\dots +\frac{1}{2^{n-1}}=\frac{1}{2^n}u_0+2-\frac{1}{2^{n-1}}.$$
Donc, $(u_n)$ est majorée car $\displaystyle \frac{1}{2^{n-1}}\underset{+\infty}{\longrightarrow}0$. \'Etant à la fois croissante est majorée, la suite $(u_n)$ converge vers $\ell \in \Rr$ et, par passage à la limite, $\displaystyle 0\leq \ell\leq 2$. Par ailleurs, l'hypothèse faite sur $u_n$ donne
$$0\leq \ell \leq \frac{\ell}{2} \Rightarrow \ell =0$$
et comme $(u_n)_{n\geq 1}$ est croissante positive, $u_n=0$ pour tout $n\geq 1$.
\end{explanations}
\end{question}
}
